% 2003 年粤桂卷第 20 题：海滨城市与台风路径示意图
\begin{tikzpicture}[>=stealth,scale=0.85]
  % 东北方向标
  \draw[->] (1.5, 1.2) -- (1.5, 1.7) node[above,scale=0.8] {北};
  \draw[->] (1.5, 1.2) -- (2.0, 1.2) node[right,scale=0.8] {东};

  % 坐标轴与原点
  \draw[->] (-1.4, 0) -- (1.4, 0) node[below,scale=0.8] {$x$};
  \draw[->] (0, -1.0) -- (0, 1.7) node[left,scale=0.8] {$y$};
  \node[above left,scale=0.8] at (0, 0) {$O$};

  % 海岸线
  \draw (-2.5, -1.1) to[out=15, in=230] (0,0) to[out=50, in=245] (0.9, 1.7);
  \node[above,scale=0.8] at (-0.7, -0.4) {海};
  \node[above,scale=0.8] at (-1.1, -0.6) {岸};
  \node[above,scale=0.8] at (-1.5, -0.8) {线};

  % 夹角 theta
  \draw (0.28, 0) arc (0:-80:0.28);
  \node[right,scale=0.8] at (0.32, -0.22) {$\theta$};

  % 台风位置 P 与 OP 连线
  \coordinate (O) at (0, 0);
  \coordinate (P) at (0.9, -4.5);
  \coordinate (Pprime) at (-0.1, -3.5);

  \draw (O) -- (P);

  % 初始台风区域（虚线圆）
  \draw[dashed] (P) circle (1.2);
  \node[below right,scale=0.8] at (P) {$P$};

  % 45度角与移动路径
  \draw (P) -- ++(-1.8, 0);
  \draw (P) -- (Pprime);
  \draw ([shift={(180:0.5)}]P) arc (180:135:0.5);
  \node[above left,scale=0.7] at ([shift={(160:0.55)}]P) {$45^\circ$};

  % t 小时后台风区域（实线圆与半径 r(t)）
  \draw (Pprime) circle (1.65);
  \node[above,scale=0.8] at (Pprime) {$P'$};
  \draw (Pprime) -- ++(-1.65, 0);
  \node[below,scale=0.8] at ([shift={(-0.85, 0)}]Pprime) {$r(t)$};
\end{tikzpicture}
